% =====================================================================
% PRX Quantum supplement — Regime-dependent quantum PINN benchmarking
% =====================================================================
% Compile: pdflatex supplement && pdflatex supplement
% Or:      bash paper/build_supplement.sh
%
% Every numerical claim here is verified by
% scripts/verify_paper_integrity.py against committed JSONs.
% =====================================================================
\documentclass[%
  prx,
  amsmath,
  amssymb,
  aps,
  reprint,
  superscriptaddress,
  floatfix,
]{revtex4-2}

\usepackage{graphicx}
\usepackage{booktabs}
\usepackage{hyperref}
\usepackage{xcolor}
\usepackage{siunitx}
\usepackage{microtype}
\hypersetup{colorlinks=true,linkcolor=blue,citecolor=blue,urlcolor=blue}

% Mirror the custom commands used in main.tex.
\newcommand{\relltwo}{\mathrm{relL}^2}
\newcommand{\dsmooth}{\Delta_{\mathrm{smooth}}}
\newcommand{\dbroad}{\Delta_{\mathrm{broad}}}
\newcommand{\ddiff}{\Delta_{\mathrm{diff}}}
\newcommand{\fhn}{FitzHugh--Nagumo}
\newcommand{\lv}{Lotka--Volterra}
\newcommand{\vdp}{Van der Pol}
\newcommand{\ac}{Allen--Cahn}

\begin{document}

\title{Supplementary Material:\\
  Regime-dependent advantage in quantum physics-informed neural
  networks: a pre-registered comparison on an ODE/PDE hardness
  ladder}

\author{Shawn Gibford}
\email{gibfords@gmail.com}
\affiliation{Independent}

\date{\today}

\maketitle

\tableofcontents

% =====================================================================
\section{Circuit specifications and faithfulness manifest}
\label{sup:circuit-specs}

The four quantum-PINN ansatz families used as solver/forecaster
encoders in the main text are each implemented against published
specifications; the binding manifest is
\texttt{refs/CIRCUIT\_SPECS.md} in the open-source release.
Table~\ref{sup:tab:circuits} summarizes the per-family
configuration used throughout the paper.

\begin{table*}[t]
\centering
\caption{Quantum-PINN ansatz roster. Per-scalar-circuit parameter
counts. Ansätze marked ``2D port'' add a split-qubit coordinate
encoding to the same base circuit (P3.9 design).}
\label{sup:tab:circuits}
\small
\begin{tabular}{lllrrl}
\toprule
Family & Source & Cell paradigm & PQC params & Classical params & Notes \\
\midrule
\texttt{data\_reuploading} & P\'erez-Salinas 2020 & encoder + LiquidQuantumCell &
  60 & 0 & 4 qubits, 5 layers (ODE) \\
\texttt{hardware\_efficient} & registry default & encoder + LiquidQuantumCell &
  60 & 0 & ring entanglement \\
\texttt{strongly\_entangling} & PennyLane native & encoder + LiquidQuantumCell &
  60 & 0 & alias of data\_reuploading in registry \\
\texttt{brickwall} & registry default & encoder + LiquidQuantumCell &
  60 & 0 & adjacent-pair CZ pattern \\
\texttt{te\_qpinn\_fnn} & Berger 2025 & FNN encoder + 4-qubit PQC &
  60 & 100 & 2D port: \texttt{te\_qpinn\_fnn\_2d} \\
\texttt{te\_qpinn\_qnn} & Berger 2025 & QNN encoder + 4-qubit PQC &
  84 & 0 & zero classical params; 2D port available \\
\texttt{qcpinn} & Farea 2025 & Classical pre/post + 4-qubit PQC &
  15 & 706 & largest classical overhead; 2D port available \\
\texttt{chebyshev\_dqc\_2d} & P3.7 design & 2D Chebyshev tower + HEA &
  -- & 0 & solver-task only, $n_t = n_x = 4$ qubits \\
\texttt{rf\_qrc} & Ahmed 2024 & Frozen reservoir + ridge readout &
  0 (trainable) & varies & non-liquid by construction (fixed leak\_rate) \\
\texttt{non\_liquid\_<ansatz>} & this work (P7.11) & encoder + NonLiquidQuantumCell &
  same as base $-Q$ & same & $\tau$-ablated; 4 variants \\
\bottomrule
\end{tabular}
\end{table*}

\subsection{Cell-level dynamics, all variants}

The hidden-state dynamics inside the Diffrax integration step (for
the forecaster task; the solver task does not integrate a hidden
state) are:

\begin{align}
\text{LiquidQuantumCell:}\quad
  \frac{dh}{dt} &= -\!\left(\frac{1}{\tau} + q(x)\right) \odot h
                 + A \odot q(x) \\[2pt]
\text{NonLiquidQuantumCell:}\quad
  \frac{dh}{dt} &= -q(x) \odot h + A \odot q(x) \\[2pt]
\text{ClassicalLTCCell:}\quad
  \frac{dh}{dt} &= -\frac{1}{\tau} \odot h
                 + \mathrm{MLP}([h, x]) \\[2pt]
\text{PlainNeuralODECell:}\quad
  \frac{dh}{dt} &= \mathrm{MLP}([h, x])
\end{align}

The two pairs (LQC vs NLQC) and (LTC vs PlainNeuralODE) differ by
exactly the $-(1/\tau) \odot h$ term. The cell-level math-identity
test (\texttt{tests/qlnn\_/test\_non\_liquid\_quantum\_cell.py})
verifies this identity numerically at random probe points.

% =====================================================================
\section{Per-cell error tables, full precision}
\label{sup:per-cell}

\subsection{Solver-task ODE matrix ($4 \times 4 \times 3 = 48$ cells)}
\label{sup:solver-ode}

Table~\ref{sup:tab:solver-ode} reports relative-$L^2$ error per
(family, system, seed) at full 4-decimal precision. Best-per-row
quantum family is bolded.

\begin{table*}[t]
\centering
\caption{Solver-task ODE per-cell relative-$L^2$ error
(PDEBench / FNO convention). Each row corresponds to one
(system, ansatz, qubits, depth) configuration evaluated at a
single random seed; each column reports one of the four
quantum families plus the classical PINN baseline at the
matched training budget. Bold marks the per-row best quantum
entry.}
\label{sup:tab:solver-ode}
\footnotesize
\setlength{\tabcolsep}{4pt}
\begin{tabular}{lc|cccc|c}
\toprule
System & Seed & cheb & fnn & qnn & qcp & cPINN \\
\midrule
\lv{} & 0 & 0.0996 & 0.0879 & 0.5242 & \textbf{0.0074} & 0.0117 \\
\lv{} & 1 & 0.1161 & 0.1407 & 0.5229 & \textbf{0.0029} & 0.0078 \\
\lv{} & 2 & 0.1024 & 0.1409 & 0.5240 & \textbf{0.0070} & 0.0071 \\
\vdp{} & 0 & 0.9064 & \textbf{0.8213} & 1.1796 & 3.4025 & 1.1569 \\
\vdp{} & 1 & 1.1516 & 0.8652 & \textbf{0.7622} & 0.8674 & 0.9434 \\
\vdp{} & 2 & 0.9087 & \textbf{0.8176} & 1.1890 & 2.6739 & 1.0550 \\
Lorenz & 0 & 1.0013 & 0.9950 & \textbf{0.9794} & 0.9957 & 0.9958 \\
Lorenz & 1 & 0.9994 & 0.9947 & \textbf{0.9769} & 0.9957 & 0.9957 \\
Lorenz & 2 & 0.9966 & 0.9946 & \textbf{0.9773} & 0.9943 & 0.9957 \\
\fhn{} & 0 & 0.5311 & 0.4521 & \textbf{0.3270} & 1.2126 & 0.6398 \\
\fhn{} & 1 & 0.6152 & 0.3903 & \textbf{0.3560} & 1.5467 & 0.6410 \\
\fhn{} & 2 & 0.6165 & 0.6248 & \textbf{0.3290} & 1.6251 & 1.3265 \\
\bottomrule
\end{tabular}
\end{table*}

\subsection{Solver-task PDE matrix ($4 \times 4 \times 3 = 48$ cells)}
\label{sup:solver-pde}

\begin{table*}[t]
\centering
\caption{Solver-task PDE per-cell relative-$L^2$ error
(PDEBench / FNO convention). Each row corresponds to one
(system, ansatz, qubits, depth) configuration evaluated at a
single random seed; each column reports one of the four
2D-port quantum families plus the classical PINN baseline at
the matched training budget. Bold marks the per-row best
quantum entry; ``---'' marks runs that diverged with no finite
$L^2$ measurement.}
\label{sup:tab:solver-pde}
\footnotesize
\setlength{\tabcolsep}{4pt}
\begin{tabular}{lc|cccc|c}
\toprule
PDE & Seed & cheb\_2d & fnn\_2d & qnn\_2d & qcp\_2d & cPINN \\
\midrule
Heat & 0 & 0.0553 & 0.0023 & 0.0455 & \textbf{0.0010} & 0.0054 \\
Heat & 1 & 0.0560 & 0.0031 & 0.0455 & \textbf{0.0020} & 0.0046 \\
Heat & 2 & 0.0565 & \textbf{0.0013} & 0.0458 & 0.0020 & 0.0036 \\
Burgers smooth & 0 & 0.3475 & 0.0146 & 0.3647 & \textbf{0.0112} & 0.0428 \\
Burgers smooth & 1 & 0.3522 & \textbf{0.0126} & 0.3543 & 0.0160 & 0.0191 \\
Burgers smooth & 2 & 0.3742 & 0.0246 & 0.3556 & \textbf{0.0206} & 0.0187 \\
\ac{} & 0 & 0.5783 & 0.5514 & \textbf{0.0528} & 0.8956 & 0.2018 \\
\ac{} & 1 & 0.5200 & 0.5345 & \textbf{0.0489} & 0.5620 & 0.0633 \\
\ac{} & 2 & 0.6146 & 0.0927 & \textbf{0.0534} & 0.4936 & 0.0514 \\
Burgers shock & 0 & 0.4282 & 0.2507 & 0.4268 & \textbf{0.2408} & 0.1682 \\
Burgers shock & 1 & 0.4385 & 0.1589 & 0.4227 & \textbf{0.1526} & 0.2730 \\
Burgers shock & 2 & 0.4168 & \textbf{0.1630} & 0.4473 & 0.2553 & 0.1270 \\
\bottomrule
\end{tabular}
\end{table*}

\subsection{Forecaster-task complete 2$\times$2 matrix ($n=9$)}
\label{sup:forecaster}

The body's Table~\ref{tab:forecaster-matrix} reports the all-to-all
5-quantum-family $\times$ 4-classical-baseline matrix. Here we add
the four non-liquid quantum forecaster columns (P7.11) so the full
2$\times$2 decomposition is visible at the per-cell level.

\begin{table*}[t]
\centering
\caption{Forecaster-task non-liquid quantum cells (P7.11).
Same training budget and adapter as the liquid forecasters in
the body's Table~\ref{tab:forecaster-matrix}; only the cell's
$\tau$-leak removed.}
\label{sup:tab:forecaster-nonliquid}
\footnotesize
\begin{tabular}{lc|cccc}
\toprule
System & Seed & nl\_dataru & nl\_hweff & nl\_stngent & nl\_brick \\
\midrule
\lv{} & 0 & 0.5830 & 0.6455 & 0.5830 & 0.7732 \\
\lv{} & 1 & 0.7951 & 0.7196 & 0.7951 & 0.7057 \\
\lv{} & 2 & 0.6482 & 0.6791 & 0.6482 & 0.3088 \\
\vdp{} & 0 & 1.4067 & 1.2307 & 1.4067 & --- \\
\vdp{} & 1 & 9.6720 & --- & 9.6720 & --- \\
\vdp{} & 2 & 1.2035 & --- & 1.2035 & --- \\
Lorenz & 0 & 1.2807 & 1.2699 & 1.2807 & 1.3128 \\
Lorenz & 1 & 0.9913 & 0.5100 & 0.9913 & 0.5025 \\
Lorenz & 2 & 1.0768 & 1.0544 & 1.0768 & 1.1571 \\
\bottomrule
\end{tabular}\\[2pt]
\footnotesize Note: \texttt{nl\_stngent} reproduces
\texttt{nl\_dataru} exactly because both dispatch the same
registered ansatz under the hood (pre-existing behavior of the
ansatz registry; see \texttt{src/qlnn\_/circuits/protocol.py}).
``---'' marks cells where rollout diverged beyond a stable
relative-$L^2$ measurement.
\end{table*}

% =====================================================================
\section{Bootstrap diagnostics and guard reports}
\label{sup:bootstrap}

All paired-bootstrap intervals are computed via
\texttt{src/quantum\_liquid\_neuralode/evaluation/h1\_verdict.py}
with $n_\mathrm{iter} = 10{,}000$, $\alpha = 0.05$ (95\% percentile
CI), and a fixed seed for reproducibility. The skyline guard uses
threshold $0.5$ (amendment A1) and excludes cells where the
known-RHS least-squares fit cannot reach the threshold.

\subsection{Per-bin sample counts (post-guard)}

Table~\ref{sup:tab:bins} reports the kept-sample sizes per regime
bin for each headline verdict.

\begin{table}[h]
\centering
\caption{Post-guard sample sizes per regime bin. The combined
n=24 verdict (PRIMARY SOLVER) is the largest sample in the paper.}
\label{sup:tab:bins}
\small
\begin{tabular}{lrr}
\toprule
Verdict & $n_\mathrm{smooth}$ & $n_\mathrm{broad}$ \\
\midrule
P5 forecaster (4-fam) & 6 & 3 \\
P7.5 solver raw n=9 & 6 & 3 \\
P7.6 combined n=18 & 12 & 6 \\
P7.6 HPO-best n=9 & 6 & 3 \\
P7.6 PDE-only n=9 & 6 & 3 \\
P7.8 full ladder n=24 & 12 & 12 \\
P7.10 forecaster combined (5-fam) & 6 & 3 \\
P7.11 all decomposition verdicts & 6 & 3 \\
\bottomrule
\end{tabular}
\end{table}

\subsection{Algebraic identities, exact}

Two per-cell identities, verified by the integrity script with
tolerance $10^{-3}$:

\begin{align}
\Delta_\mathrm{combined}
  &= \Delta_\mathrm{quantum\_via\_LTC} + \Delta_\mathrm{liquid\_via\_classical}
  \notag \\
  &= -0.6160 + 0.1153 = -0.5007 \checkmark \\[4pt]
\Delta_\mathrm{combined}
  &= \Delta_\mathrm{liquid\_via\_quantum} + \Delta_\mathrm{quantum\_via\_nonliquid}
  \notag \\
  &= -0.3339 + (-0.1668) = -0.5007 \checkmark
\end{align}

Both identities are derived from the same per-cell records; only
the slot-assignment in the bootstrap differs. The exactness
($|\text{diff}| < 10^{-3}$) is the cleanest verification that the
decomposition is a real mathematical structure, not a sampling
artifact.

\subsection{Q-Q residual diagnostics (post-hoc; see Amendment A14)}
\label{sup:qq-diagnostics}

The diagnostic figure suite includes a residual-distribution panel
(\texttt{fig\_residual\_analysis}, middle column) whose original
source comment mis-labelled it as a Q-Q plot. The implementation is
actually a histogram with a fitted Normal PDF overlay --- not a
quantile-quantile analysis. Because the Normal overlay was visually
misleading (both stacks' residuals are clearly non-Normal), P8 polish
adds two reusable, dataset-agnostic helpers in
\texttt{scripts/make\_diagnostic\_figures.py} ---
\texttt{\_qq\_panel\_vs\_normal} and \texttt{\_two\_sample\_qq} ---
plus a new figure \texttt{fig\_qq\_analysis} that applies them to the
canonical archived comparison at $h=3$ (\texttt{classical\_H4}
hidden-size 4 versus \texttt{QLNN\_4q3L}, seed-mean residuals on the
held-out test windows). Both helpers take any
\texttt{(residuals, label)} pair, so the same panels can be reapplied
to post-pivot solver or forecaster residuals with no code
duplication.

Table~\ref{sup:tab:qq} reports the per-stack Shapiro--Wilk normality
test and the two-sample Kolmogorov--Smirnov test between the two
empirical residual distributions.

\begin{table}[h]
\centering
\caption{Q-Q residual diagnostics for the canonical archived
comparison at $h=3$ (seed-mean test-window residuals; see
\texttt{paper/figures/fig\_qq\_analysis.pdf}). Both per-stack
distributions are clearly non-Normal; the two-sample Kolmogorov--
Smirnov test cannot reject equal distributions, quantifying the
visually-suggested ``the two histograms look identical'' observation
and supporting the project's archived ``both near persistence at
$h=3$'' verdict.}
\label{sup:tab:qq}
\small
\begin{tabular}{lccl}
\toprule
Test & Classical $H{=}4$ & QLNN $4q3L$ & Interpretation \\
\midrule
Shapiro--Wilk $W$ & 0.892 & 0.885 & both non-Normal \\
Shapiro--Wilk $p$ & $1.64{\times}10^{-5}$ & $8.91{\times}10^{-6}$ & reject $H_0{:}$ Normal \\
Two-sample KS $D$ & \multicolumn{2}{c}{0.085} & --- \\
Two-sample KS $p$ & \multicolumn{2}{c}{0.964} & cannot reject equal dist. \\
\bottomrule
\end{tabular}
\end{table}

The harness is committed and unit-tested (T1 registry assertion
updated $7 \to 8$ callables). No headline number is touched and
\texttt{scripts/verify\_paper\_integrity.py} continues to exit $0$.

% =====================================================================
\section{Reproduction and audit}
\label{sup:repro}

All scripts and data referenced below ship with the open-source
release at \url{https://github.com/shawngibford/qlnn}; the
specific commit/tag anchors are listed in \S\ref{sup:repro}
under ``GitHub release tags''.

\subsection{One-command reproduction}

The entire benchmark regenerates from scratch via:
\begin{verbatim}
bash scripts/reproduce_paper.sh
\end{verbatim}
Estimated wall-clock on a single MacBook Pro M1: approximately
8\,hours. Each sweep is independent and the script is structured
so individual phases (P3.5--P7.11) can be re-run in isolation.

\subsection{Mechanical integrity gate}

\begin{verbatim}
PYTHONPATH=src .venv/bin/python \
  scripts/verify_paper_integrity.py
\end{verbatim}
The script reads committed JSON files from \texttt{results/} and
matches each cited numerical claim in the paper body against its
on-disk value with explicit tolerances. Returns exit-0 if and only
if every claim verifies. At the submission commit the script gates
16+ $H_1$ outcomes (5 solver-task sensitivity points + 5
forecaster-task decomposition verdicts + 3 PDE-only and
ODE-only sub-verdicts) plus the two algebraic identities and the
H3 mechanism trend.

\subsection{GitHub release tags}

\begin{itemize}
  \item \texttt{v0.1-prepaper-rigor} — post-P7.6 (HPO symmetry + 16
        gated outcomes). Anchored at commit \texttt{2309faa}.
  \item \texttt{v0.2-airtight-matrix} — post-P7.8 (full ladder
        n=24 PRIMARY SOLVER verdict + Kuramoto/KdV deferral
        documentation). Anchored at commit \texttt{79cf435}.
  \item \texttt{v1.0-paper-draft} — submission anchor (to be
        created at acceptance / preprint posting).
\end{itemize}

Repository: \url{https://github.com/shawngibford/qlnn}.

% =====================================================================
\section{Pre-registration and amendments}
\label{sup:prereg}

The full pre-registration document is at
\texttt{ODE\_PDE\_PRE\_REG.md} in the repository (locked
2026-05-19 before any benchmark results were generated;
SHA stamped). Section~7's decision rule is applied mechanically by
\texttt{verify\_paper\_integrity.py}.

The complete amendment record (A1--A13) is at
\texttt{PRE\_REG\_AMENDMENT.md}. Table~\ref{sup:tab:amendments}
summarizes the headline status of each amendment.

\begin{table*}[t]
\centering
\caption{Pre-registration amendments (A1--A13). Status column
flags PRIMARY = canonical headline verdict; SUPERSEDED = the
amendment is documented but a later amendment takes precedence
for the paper's narrative; DISCLOSED = methodological choice
made open without superseding. Note: A14 (Q-Q residual
diagnostics, see \S\ref{sup:qq-diagnostics}) is a post-hoc
amendment introduced after the pre-registration set was
locked, and is intentionally not part of the pre-registered
A1--A13 record; it is documented in
\texttt{PRE\_REG\_AMENDMENT.md}.}
\label{sup:tab:amendments}
\small
\begin{tabular}{lll}
\toprule
\# & Subject & Status \\
\midrule
A1 & skyline\_threshold = 0.5 & LOCKED \\
A2 & $n=3$ seeds; 3-then-4 ODE systems & DISCLOSED \\
A3 & fixed hyperparameters $\to$ sensitivity sweep added at P7.6 & DISCLOSED + sensitivity \\
A4 & MLP capacity $3.3\times$ violates §6 factor-of-2 (conservative direction) & DISCLOSED \\
A5 & \vdp{} at $\mu=5$ near learnability boundary & DISCLOSED \\
A6 & underfit guard active on solver only & DISCLOSED \\
A7 & strict-vs-raw verdict reporting at P7.5 & DISCLOSED \\
A8 & H3 trend $\rho = +0.518$ not significant at $n=9$; LOO-robust & DISCLOSED \\
A9 & symmetric QLNN HPO (P7.6 c1) & SUPERSEDED-BY-A11 \\
A10 & combined ODE+PDE n=18 (P7.6 c2) & SUPERSEDED-BY-A11 \\
A11 & full-ladder n=24 sign-flip (P7.8) & \textbf{PRIMARY SOLVER} \\
A12 & classical LTC + 3-verdict forecaster decomposition (P7.10) & PRIMARY FORECASTER (via classical path) \\
A13 & non-liquid QLNN + complete 2$\times$2 + $\tau$-cross-check disagreement (P7.11) & \textbf{PRIMARY FORECASTER (complete)} \\
\bottomrule
\end{tabular}
\end{table*}

% =====================================================================
\section{Deferred follow-up work}
\label{sup:deferred}

Two pre-registered systems and three optimization techniques are
deferred to the follow-up paper. Each is documented here with
its current state (where applicable, the mechanism gate has been
cleared so the follow-up can proceed without architectural risk).

\paragraph{Kuramoto (12-dimensional coupled oscillators).}
Per-component scalar circuits scale linearly with state dimension.
At our default config (4 qubits per component, $L=1$
reuploading), a 12D system projects to $\approx 7$\,hr per cell.
The per-component dispatcher works correctly (single-cell smoke
verified); the full sweep is the follow-up paper's first compute
item.

\paragraph{KdV (Korteweg--de Vries, $u_t + 6 u u_x + u_{xxx} = 0$).}
\label{sup:kdv-gate}
Third-order spatial derivative requires triple-nested
reverse-mode autodiff (\texttt{jacrev}$^3$) through the QNode.
The mechanism gate at \texttt{scripts/run\_p7\_8\_kdv\_gate.py}
PASSED: $u_{xxx}$ values are finite and non-trivial, JIT compile
takes 4.55\,s, and the per-point cost ratio is $0.43\times$ the
baseline second-order autodiff (XLA fuses the third derivative
tightly). The integrated training cost remains prohibitive at
$\approx 12$\,s per gradient step, projecting $\approx 8$\,hr per
seed across the four quantum families + classical PINN baseline.

\paragraph{Quantum Natural Gradient~\cite{stokes2020quantum}.}
Targets the quantum-circuit training-inefficiency identified by
the LTC decomposition's $\Delta_\mathrm{quantum\_via\_LTC} =
-0.616$. Uses the Quantum Fisher Information matrix as the
optimizer preconditioner.

\paragraph{Causal training weighting (Wang, Sankaran,
Perdikaris)~\cite{wang2022causality}.}
Weights collocation points by inverse time-order so that residuals
at early times converge first. Targets the \vdp{} stiff-dynamics
failure mode (where multiple cells diverge to
$\relltwo > 1$ at the program's compute budget).

\paragraph{Higher data-reuploading depth.}
The H3 analysis found Fourier $K_{\max} = 3$ at $L = 1$ across
all four ansätze: at the theoretical ceiling
$K_{\max} = L \cdot n$ with $L=1, n=4$. The
Schuld--Sweke--Meyer relation predicts a 5-fold increase
in accessible Fourier bandwidth at $L = 5$. This is the most
direct test of whether the LTC decomposition's
$\Delta_\mathrm{liquid\_via\_classical} = +0.115$ smooth-bin
inductive bias can be recovered by giving the quantum side more
expressivity.

% =====================================================================
\bibliographystyle{apsrev4-2}
\bibliography{references}

\end{document}
